\chapter{Project overview}
\label{ch:ch1}
\section{Problem and objective}
A compromised container can execute unexpected processes, open outbound connections
or access credentials available to its workload. Detecting that behaviour is only
one part of the response: deleting the pod immediately can interrupt service and
destroy volatile evidence, while leaving it connected permits further activity.

This project implements a controlled-environment prototype that connects runtime
alerts to Kubernetes API actions. Its objective is to detach a selected compromised
pod from normal service, retain it for investigation and allow the workload controller
to create a replacement. Cilium policies provide the configured network-quarantine
mechanism. The response is asynchronous; changing a label is not proof that the
network policy has already been enforced.

\section{Architecture}
The implementation separates event detection, event forwarding, response decisions
and network enforcement:
\begin{enumerate}
\item \textbf{Falco} evaluates syscall events against runtime-security rules.
\item \textbf{Falcosidekick} receives alerts over HTTP and forwards selected events
to the response webhook with a shared authentication token.
\item \textbf{The Python handler} validates event scope and updates the target pod
through the Kubernetes API, using conflict-aware metadata patches.
\item \textbf{Kubernetes controllers} reconcile Service membership and replacement
replicas after the workload selector labels are removed.
\item \textbf{Cilium} applies explicit ingress and egress deny rules to endpoints
selected by the quarantine label.
\end{enumerate}
The target application, \texttt{vulnapp}, is a minimal Flask network-diagnostic form
with a deliberate command-injection vulnerability. Its source, container definition,
deployment manifests and a Linux ptrace demonstration are included in the repository.

\section{Scope and trust boundary}
The prototype is intended for an isolated, disposable environment. The vulnerable
application must not be exposed to public or untrusted networks. The handler is
restricted to the configured namespace and workload label; cluster administrators,
node access and users who can create pods or modify labels remain trusted.

The project is not a production incident-response platform. It does not acquire
memory, verify exported evidence, implement a durable event queue or guarantee
continuous service. Two application replicas reduce exposure to a single-pod
incident but do not establish zero downtime. Network enforcement, controller
reconciliation and failure recovery must be measured in the deployment environment.

\section{Implementation and verification status}
The response handler, deployment hardening, network-policy manifests and automated
tests are implemented. Sixteen regression tests passed in the Python 3.11 handler
image under non-root, read-only and network-disabled execution. Both container
images were built, Helm output was inspected, and standard Kubernetes manifests
passed server-side dry-run validation.

\textbf{End-to-end Falco-to-Cilium containment remains unverified for this revision.}
The available validation environment did not include Cilium. Detection scenarios
in Chapter~\ref{ch:ch3} specify reproducible procedures and acceptance criteria,
not a claim that the complete chain has already passed those checks.
